% Superseded by the mass argument on October 23, 2026.
\paragraph{Idea of the proof.}
By Corollary~\ref{cor:one-point} it is enough to find one point of strict
convexity between $m^*$ and $2m^*+1$. If there were none, $P_N$ would be
concave on $[m^*,2m^*+2]$, and then its value halfway would be at least half of
its maximum. The Landau limit shows that the value halfway is much smaller.

\begin{proof}
Let $N$ be large, and let $\log N\le x\le2\log N$. Then $\varepsilon\le2\log N/N$, so
\eqref{eq:H} holds and $\varepsilon<\frac12$. Also $(1+\log x)/x\le(1+\log\log N)/\log N$,
$x/N\le2\log N/N$ and $x\delta_B=O(x^3\log^2N/N)=O(\log^5N/N)$. So by
Theorem~\ref{thm:landau}(b) there are numbers $\Theta_N\to0$, which depend on $N$ alone, such
that for every $x$ in this range
\begin{equation}\label{eq:llt-used}
 \bigl\lvert xf_N(m)-f_L(\lambda_m)\bigr\rvert\le\Theta_N\qquad\text{for every integer }m\ge0 .
\end{equation}
The density $f_L$ is continuous, because it is the Fourier transform of an
integrable function. It is unimodal with mode $\lambda_{\mathrm L}$
(Lemma~\ref{lem:landau-mode}), so it is nonincreasing on
$[\lambda_{\mathrm L},\infty)$, and it is integrable. So
$f_L(\lambda)\to0$ as $\lambda\to\infty$. We write $o(1)$ for a quantity that
is bounded in absolute value by a function of $N$ alone that tends to $0$. So every bound
below holds uniformly in $x$.

\emph{Step 1. The mirror term.} Let $0\le m\le N/2$. Then
$\lambda_{N-m}\ge\frac N{2x}-\log x\ge\frac N{4\log N}-\log(2\log N)$, which tends to
$\infty$. So \eqref{eq:llt-used} gives
$xf_N(N-m)\le f_L(\lambda_{N-m})+\Theta_N=o(1)$. Since
$P_N(m)=P_N(N-m)=\frac12[f_N(m)+f_N(N-m)]$, we get
\[
 2xP_N(m)=f_L(\lambda_m)+o(1)\qquad\text{uniformly for }0\le m\le N/2 .
\]

\emph{Step 2. The mode.} Let $m_0$ be the integer part of
$x\log x+\lambda_{\mathrm L}x$. Then $0\le m_0\le N/2$ and
$\lvert\lambda_{m_0}-\lambda_{\mathrm L}\rvert\le1/x$, so Step~1 and the continuity of
$f_L$ give $2xP_N(m_0)=f_L(\lambda_{\mathrm L})+o(1)$. Now $P_N(m^*)\ge P_N(m_0)$ and
$m^*\le N/2$, so Step~1 at $m^*$ gives
$f_L(\lambda_{m^*})\ge f_L(\lambda_{\mathrm L})-o(1)$. By
Lemma~\ref{lem:landau-mode}(c) this forces
$\lvert\lambda_{m^*}-\lambda_{\mathrm L}\rvert\le1$ for large $N$. Hence
\[
 x\log x+(\lambda_{\mathrm L}-1)x\le m^*\le x\log x+(\lambda_{\mathrm L}+1)x,\qquad
 2xP_N(m^*)\ge f_L(\lambda_{\mathrm L})-o(1).
\]
In particular $m^*\to\infty$ and $2m^*+2\le N/2$ for large $N$.

\emph{Step 3. Halfway.} Let $v=m^*+\lfloor m^*/2\rfloor$. Then
$v\ge\frac32m^*-1$, so by Step~2
\[
 \lambda_v\ge\tfrac12\log x+\tfrac32(\lambda_{\mathrm L}-1)-\tfrac1x
 \ge\tfrac12\log\log N+\tfrac32(\lambda_{\mathrm L}-1)-1 ,
\]
which tends to $\infty$.
Also $v\le N/2$. So Step~1 gives $2xP_N(v)=f_L(\lambda_v)+o(1)=o(1)$. Since
$f_L(\lambda_{\mathrm L})>0$, Step~2 shows that $P_N(v)<\frac12P_N(m^*)$ for large $N$.

\emph{Step 4. A point of convexity.} Suppose that $D_N(a)\le0$ for every
$m^*<a\le2m^*+1$. Then $P_N$ is concave on the integers of $[m^*,2m^*+2]$.
With $u=m^*$ and $w=2m^*+2$ this gives
\[
 P_N(v)\ge\frac{w-v}{w-u}P_N(u)+\frac{v-u}{w-u}P_N(w)\ge\frac{w-v}{w-u}P_N(m^*)
 \ge\tfrac12P_N(m^*),
\]
because $w-v\ge\frac12(m^*+2)=\frac12(w-u)$. This contradicts Step~3. So
$D_N(a_1)>0$ for some $m^*<a_1\le2m^*+1$, and $a_1\le N/2$ by Step~2.
By Theorem~\ref{thm:dip}, $P_N(1)=P_N(N-1)>\frac x2E_N>P_N(0)$, since $x\ge\log N>2$. As
$\varepsilon<\frac12$, the second case of Corollary~\ref{cor:one-point} applies and gives the
claim. Finally, Theorem~\ref{thm:erw-crossing}(b) gives $\log N\le x_N\le2\log N$ for every
large even $N$, which is the last claim.
\end{proof}

\begin{remark}[what is still open]\label{rem:convex-open}
The remaining gap concerns the prescribed convexity window. The minimum
and two-hump shape hold for every even $N\ge8$.
The threshold $N_1$ in Theorem~\ref{thm:convex-large} is not computed.
Even without its error, Step~3 needs
$f_L(\frac32\lambda_{\mathrm L}+\frac12\log x)<f_L(\lambda_{\mathrm L})/2$.
The numerical half-height point $2.432$ gives $x>252.8$, or a crossing
length above about $3.5\cdot10^{56}$. Thus that argument cannot meet the
finite checks. Corollary~\ref{cor:one-point} would close the gap from
$D_N(a_1)>0$ at one point $m^*<a_1\le2m^*+1$ (and
Proposition~S3.1(b) of Online Resource~1 gives the odd-$N$ version).
The local mass error is too large to prove this second-difference sign.
\end{remark}
